%!TEX root = ../main.tex
\chapter{Introduction}
\label{ch:introduction}

\section{Problem Description}

Autonomous navigation of a robotic platform in an unknown environment requires
solving, continuously and in real time, two problems: knowing where the robot is,
and knowing what surrounds it, so as to move towards a goal without hitting
anything. In common practice these two problems are addressed by pipelines built
around one specific platform: a given set of sensors, mounted in given positions,
on a vehicle with a given kinematics. Changing any one of these elements ---
replacing a LiDAR with a depth camera, moving from a rover to a multirotor ---
typically means rewriting significant parts of the system, in perception as well
as in control.

This thesis is set within the Vulturna project, an autonomy pipeline designed
to be \emph{hardware-independent}: perception, state estimation, mapping and
navigation are built so that the same system runs on different robots --- ground
or aerial, with different sensors, geometries and kinematic limits --- without
changing the code. Which sensors are present, how they are mounted and within
which limits the platform moves are configuration data, not knowledge embedded
in the software. Downstream, the pipeline delivers a unified representation of
the surroundings that is independent of the source that produced it.

The problem addressed by this thesis is the last link of that chain: deciding
the motion. The approach is learnt, not programmed. No explicit rule describes
how to avoid an obstacle or how to approach a goal; a neural network --- called
the \emph{microbrain} within the project --- observes the space around the robot
and directly produces the velocity command, taking exactly the place a human
pilot would take. For this to be possible on an agnostic platform, the network
must meet three requirements that rarely coexist. It must consume an observation
whose form does not depend on the sensor that produced it, and therefore work on
a sparse geometric representation rather than on an image. It must produce a
command that is valid for vehicles with different kinematics --- a rover that
cannot climb, a multirotor that translates in every direction --- by reading the
vehicle's limits from the observation itself. And it must have memory: a single
snapshot of the local surroundings is not enough to decide where to go, because
the local map forgets what the robot no longer sees, data arrive at irregular
intervals, and the correct manoeuvre depends on what has just happened.

It is this last requirement that motivates the choice of \emph{Liquid Neural
Networks} (LNN): a family of continuous-time recurrent networks, derived from
Liquid Time-Constant networks, in which the internal state evolves according to
a differential equation whose integration step is the real interval elapsed
between two observations. Compared with classical recurrent networks they offer
a memory that accounts for physical time, a small size and --- in the variants
with a sparse wiring inspired by biological neural circuits --- an interpretable
structure. They were proposed precisely for robot and vehicle control tasks, but
their application to a multi-robot, multi-sensor platform fed with point clouds
rather than images has not been explored. This thesis studies whether and how
these networks can serve as the microbrain of an agnostic platform, and which of
their variants is best suited to the task.

\section{Objectives}

The work is organised around three objectives, which also follow the
chronological order in which they were addressed.

\paragraph{The platform.}
The first objective, carried out during the internship that preceded the
thesis, was the realisation of the robotic pipeline everything else rests on: a
ROS~2 system capable of handling, in a sensor-agnostic way, an arbitrary and
varied set of sensors, and of fusing their data into a single estimate of the
robot's pose and a single representation of its surroundings, suitable for
deciding the motion. The pipeline runs identically in simulation and on real
hardware. For this thesis it is, at once, the producer of the observation the
network consumes, the recorder of the training data, and the environment in
which the network is put to the test.

\paragraph{The networks.}
The second objective is the analysis of Liquid Neural Networks, considered
fundamental to the problem for their ability to maintain a memory consistent
with physical time. The analysis starts from the mathematical formulation of
the three cells considered --- Liquid Time-Constant (LTC), Closed-form
Continuous-time (CfC) and Liquid Resistance-Capacitance Unit (LRCU) --- and of
the two wiring schemes, dense and sparse according to the Neural Circuit Policy
(NCP), and continues with their implementation in PyTorch, first as reusable
modules and then as complete networks. Around the cells sit a sparse
convolutional encoder, which reduces the voxel cloud produced by the pipeline
to a compact descriptor, and the auxiliary modules that make the network usable
on a heterogeneous fleet of robots: observation normalisation, the robot-model
embedding and output calibration. The modules and the networks are validated
first on synthetic tasks and then on the recorded data.

\paragraph{Training and comparison.}
The third objective is the design of the training of these networks and their
comparison on the problem, in order to identify the best candidate. The strategy
has two stages. In the first, the networks learn by imitation from a corpus of
flights recorded on the pipeline, in which every sample pairs what the robot
observed with the command the pilot gave; this stage produces the weights the
robot flies with and those the next stage starts from. In the second, the
networks are refined by reinforcement learning, flying in a fleet of parallel
simulations and receiving a reward designed so that reaching the goal, without
losing the vehicle, is numerically the best policy. Both stages are instrumented
with a set of metrics designed to distinguish genuine learning from its
imitations, and are validated on simpler control tasks before being applied to
the problem. The comparison between the architectures is carried out on those
metrics.

\section{Contributions}

\section{Thesis Outline}

The rest of the thesis is organised as follows.

Chapter~\ref{ch:related} places the work with respect to the literature:
learnt navigation for aerial and ground robots, Liquid Neural Networks and
their applications to control, point-cloud encoders, and the combination of
imitation and reinforcement learning.

Chapter~\ref{ch:ros2} presents the ROS~2 pipeline built during the internship:
what it does, the hardware-agnostic principle it is built on, how data from
heterogeneous sensors are fused, and the observation it delivers to the network
--- the observation space on which everything that follows is built.

Chapter~\ref{ch:lnn} presents the Liquid Neural Networks: the formulation of
the LTC, CfC and LRCU cells and of the NCP wiring, the sparse convolutional
encoder and the auxiliary modules, their implementation in PyTorch and the
networks that result.

Chapter~\ref{ch:metrics} defines, before any result, the metrics through which
the two training stages are read: what each one measures, why it was
introduced, and how they are to be read together. It is a self-contained
chapter, meant as a reference for the chapters that follow.

Chapter~\ref{ch:training} describes the training: the two-stage strategy; for
the supervised stage, the dataset and its pipeline, the loss function, the
output calibration and the results; for the reinforcement stage, the algorithm
(a recurrent variant of PPO with GAE), the environment and the episode policy,
the reward design, the fine-tuning from the imitated weights, the validation on
a standard control task and the results on the fleet. The chapter closes with
the comparison between the architectures.

Chapter~\ref{ch:conclusions} summarises the results obtained, discusses the
limits of the work and points to future developments.
